% Bilaga A, satt ur info/simulation_workflow.md.

\chapter{Simulering med GHDL}
\label{app:sim}

Det här är bokens permanenta referens för simulering. Varje kapitel som verifierar en modul mot en
testbänk pekar hit i stället för att upprepa stegen. \Secref{c2:sec:testbenches} är den mjukare
introduktionen, skriven för första gången du kör en testbänk över huvud taget; den här bilagan är den
du kommer tillbaka till, och den täcker grupprojektets bygge i synnerhet.

% ------------------------------------------------------------------------------------------
\section{Varför simulera}
\label{app:sim:why}

\Chapref{c1:ch} till \ref{c8:ch} verifierar varje konstruktion för hand, i CircuitVerse eller direkt
på DE0-CV. Det fungerar för små kretsar med en handfull ingångar som en människa kan växla.

Det slutar fungera så snart en konstruktions korrekthet beror på exakta tidsförhållanden över flera
cykler mellan flera moduler, vilket är grupprojektets situation rakt igenom.

\begin{itemize}
  \item En \textbf{testbänk} är ett litet, utdelat VHDL-program som driver en konstruktions ingångar
    och kontrollerar dess utgångar automatiskt.
  \item \textbf{GHDL} är verktyget som kompilerar och kör den, helt i programvara, utan att någon
    FPGA är inblandad.
  \item Du ombeds aldrig skriva den testbänk som \emph{verifierar} en modul du bygger. Varje sådan
    är utdelad; du kör den och läser vad den rapporterar.
  \item Några övningar bjuder in dig att skriva en kort \emph{provisorisk} rigg, tjugo rader som
    matar in din egen bitsekvens i en modul och \code{report}:ar vad som kommer ut. Det är
    utforskande, och alltid frivilligt till formen: ingenting du bedöms på beror på en testbänk du
    själv skrivit.
\end{itemize}

% ------------------------------------------------------------------------------------------
\section{Att installera GHDL}
\label{app:sim:install}

På Ubuntu eller WSL:

\begin{shell}
sudo apt-get update && sudo apt-get install -y ghdl
\end{shell}

Bekräfta att det gick vägen:

\begin{shell}
ghdl --version
\end{shell}

GHDL 3.x eller senare. Kursens CI kör samma GHDL-kommandon som nedan, så det som passerar lokalt
passerar också i CI.

% ------------------------------------------------------------------------------------------
\section{Var filerna bor}
\label{app:sim:layout}

Grupprojektets VHDL ligger i \textbf{två platta kataloger i repots rot}: modulerna du skriver, vid
sidan av de utdelade testbänkarna (\file{*_tb.vhd}) som kontrollerar dem. Det finns inga kataloger
per pass och ingenting att kopiera mellan dem; \file{can_def.vhd} i synnerhet skrivs \emph{en} gång,
i \chapref{c10:ch}, och varje senare kontrollermodul läser den enda filen.

Uppdelningen mellan de två katalogerna är protokollgränsen: allt i \file{controller/} vet något om
CAN och ingenting om SPI, och allt i \file{bridge/} tvärtom.

\begin{textdrawing}
controller/   can_def.vhd            kap 10, din
              can_controller.vhd     kap 11, din (först bara portar, växer genom kap 17)
              meta_prev.vhd          kap 12, din
              bit_timer.vhd          kap 12, din
              crc15.vhd              kap 13, din
              tx_shift_reg.vhd       kap 14, din
              rx_shift_reg.vhd       kap 15, din
              + de sex utdelade *_tb.vhd

bridge/       register_bank.vhd      kap 18, din
              spi_reg_bridge.vhd     kap 19, din
              can_spi_node.vhd       kap 19, din
              + utdelade spi_slave.vhd, spi_def.vhd,
                register_bank_tb.vhd, spi_reg_bridge_tb.vhd
\end{textdrawing}

De två katalogerna innehåller från början bara de utdelade filerna: åtta testbänkar, plus
\file{spi_slave.vhd} och \file{spi_def.vhd}. Varje kapitel från \chapref{c10:ch} lägger till en
modul, med undantaget att \file{can_controller.vhd} skrivs tidigt (\chapref{c11:ch}) som en entitet
med en tom arkitektur och sedan \emph{växer}: \chapref{c16:ch} fyller i dess sändväg och
\chapref{c17:ch} dess mottagväg, utan att någon ny fil dyker upp för någondera. Det är därför
\code{can_controller_tb}, den enda testbänk som prövar hela kontrollern, är grindad på mer än att
filerna finns; \secref{app:sim:make} återkommer till det.

% ------------------------------------------------------------------------------------------
\section{De tre GHDL-stegen}
\label{app:sim:steps}

Varje VHDL-konstruktion (och testbänk) går igenom samma tre steg, i den här ordningen:

\begin{enumerate}
  \item \textbf{Analysera} (\code{ghdl -a}): tolkar en \file{.vhd}-fil, kontrollerar den mot
    syntax- och typfel, och registrerar dess designenheter (entiteter, arkitekturer, paket) i ett
    arbetsbibliotek. Beroenden måste analyseras före de enheter som använder dem:
    \file{can_def.vhd} före varje modul som säger \code{use work.can_def.all;}, och en modul före
    sin egen testbänk.
  \item \textbf{Elaborera} (\code{ghdl -e}): löser upp en viss toppnivåentitets fullständiga
    instansieringshierarki till något körbart.
  \item \textbf{Köra} (\code{ghdl -r}): simulerar den på riktigt.
\end{enumerate}

För en modul med utdelad testbänk, till exempel \chapref{c12:ch}:s \code{meta_prev}, som inte läser
något paket alls och därför bara behöver sina egna två filer:

\begin{shell}
cd controller
ghdl -a --std=93 meta_prev.vhd meta_prev_tb.vhd
ghdl -e --std=93 meta_prev_tb
ghdl -r --std=93 meta_prev_tb --assert-level=error
\end{shell}

Varje annan kontrollermodul läser \code{can_def}, så paketet kommer först:

\begin{shell}
cd controller
ghdl -a --std=93 can_def.vhd bit_timer.vhd bit_timer_tb.vhd
ghdl -e --std=93 bit_timer_tb
ghdl -r --std=93 bit_timer_tb --assert-level=error
\end{shell}

\file{bridge/} fungerar likadant, med en krumelur: \code{register_bank} läser \code{can_def}, som
ligger en katalog bort.

\begin{shell}
cd bridge
ghdl -a --std=93 ../controller/can_def.vhd register_bank.vhd register_bank_tb.vhd
ghdl -e --std=93 register_bank_tb
ghdl -r --std=93 register_bank_tb --assert-level=error
\end{shell}

\begin{itemize}
  \item \code{--std=93} väljer VHDL-93 genom hela kursen. Varje utdelad testbänk avslutar sig själv
    snyggt, genom att sätta en \code{done}-signal som klockgenereringsprocessen kontrollerar och
    sedan anropa \code{wait;}, i stället för att förlita sig på någon ”stoppa simulatorn”-procedur
    ur en senare standard.
  \item \code{--assert-level=error} säger åt GHDL att stanna simuleringen, med en utgångskod skild
    från noll, vid varje misslyckad kontroll som rapporterats på nivån \code{error} eller högre.
    \textbf{Skicka alltid med den när du kör en projekttestbänk.} Nästa avsnitt säger varför den inte
    är valfri här.
\end{itemize}

\subsection{Varför \code{--assert-level=error} spelar roll, och var den inte gör det}
\label{app:sim:assert}

GHDL stannar på en misslyckad assertion bara när assertionens allvarlighetsgrad når assert-nivån,
och GHDL:s förvalda nivå är \code{failure}. Bokens två halvor ligger på var sin sida om den gränsen,
vilket är värt att veta eftersom det förklarar varför flaggan ser valfri ut i del~\ref{part:vhdl} och
inte är det i del~\ref{part:can}:

\begin{booktable}{lcLL}
  & \thead{Grad} & \thead{Utan flaggan} & \thead{Med flaggan} \\
  \midrule
  Kapitlens testbänkar & \code{failure} & stannar, avslutar 1 & stannar, avslutar 1 \\
  Projektets testbänkar & \code{error} & \textbf{skriver ut, kör vidare, avslutar 0} & stannar,
    avslutar 1 \\
\end{booktable}

För en projekttestbänk gör alltså ett utelämnande av flaggan om ett verkligt fel till ett meddelande
som rullar förbi ovanför en nollställd utgångskod, vilket är precis vad ett ”tyst godkänt” ser ut som.
Skicka med den varje gång så slutar skillnaden spela roll.

\code{--stop-time=10ms} är värd att lägga till av samma slags skäl. En testbänk som väntar på en
händelse din modul aldrig producerar \emph{fallerar} inte; den hänger. Stopptiden gör om det till en
körning som tar slut. Den står inte i kommandona ovan därför att projektets testbänkar alla begränsar
sina egna väntor, men den kostar ingenting och \secref{c2:sec:testbenches} gör den till en vana.

Varje testbänk rapporterar bara godkänt eller underkänt (och, vid underkänt, orsaken), ingenting
annat. En godkänd körning slutar med en rad som denna:

\begin{console}
bit_timer_tb.vhd:148:9:@7312ns:(report note): bit_timer: all checks passed!
\end{console}

En underkänd körning stannar tidigt med ett \code{(assertion error)}-meddelande som namnger exakt
vilken kontroll som fallerade och varför. Läs det meddelandet; det är skrivet för att tala om
specifikt vad som gick fel.

% ------------------------------------------------------------------------------------------
\section{Att bygga allt på en gång}
\label{app:sim:make}

\code{make build-project}, från repots rot, gör allt det här automatiskt för varje testbänk vars
moduler du skrivit. Det är användbart för att bekräfta att hela ditt träd fortfarande passerar, men
det är förståelsen av de tre enskilda stegen ovan som låter dig felsöka en enskild modul snabbt under
ett pass.

Eftersom projektet byggs upp en modul i taget \textbf{hoppar} \code{make build-project} över varje
testbänk vars moduler ännu inte finns, i stället för att fallera på dem, och säger det:

\begin{console}
==> meta_prev_tb
controller/meta_prev_tb.vhd:114:9:@311ns:(report note): meta_prev: all checks passed!
==> bit_timer_tb (skipped: no bit_timer.vhd yet)
\end{console}

Så ser \chapref{c12:ch} ut halvvägs: \file{meta_prev.vhd} skriven och godkänd, \file{bit_timer.vhd}
kvar att göra. \file{can_def.vhd} namnges inte som saknad, eftersom \chapref{c10:ch} redan skrev den.

Överhoppad är inte underkänd, med flit. Projektet löper över tio pass och det mesta av trädet saknas
under de flesta av dem; ett bygge som lyste rött tills den sista modulen landade skulle inte säga
något användbart under någon av de nio föregående. Antalet testbänkar som faktiskt kör växer alltså
allteftersom: inga efter \chapref{c10:ch} eller \ref{c11:ch}, två efter \chapref{c12:ch}, sedan en
till i vart och ett av \chapref{c13:ch}, \ref{c14:ch} och \ref{c15:ch}, en till i \chapref{c17:ch},
en i \chapref{c18:ch}, och den sista i \chapref{c19:ch}.

\begin{aside}
\code{can_controller_tb} är \textbf{grindad på mer än att filen finns}, och det är värt att veta
innan det förbryllar dig. \file{can_controller.vhd} skrivs i \chapref{c11:ch} som en entitet med tom
arkitektur och växer genom \chapref{c16:ch} och \ref{c17:ch}, så från \chapref{c15:ch} och framåt
finns varje fil testbänken behöver medan kontrollern fortfarande inte kan ta emot en ram. Att köra
den då skulle fallera, två gånger om: en gång i \chapref{c15:ch}, och igen i \chapref{c16:ch}, som
kan sända men ännu inte ta emot. Eftersom testbänken fungerar genom att låta en nod ta emot en annans
ram är det den faktiskt kräver en \textbf{mottagväg}, och den kommer i \chapref{c17:ch}.
Byggskriptet tittar därför inuti din \file{can_controller.vhd} efter \chapref{c17:ch}:s CRC-grindning
på mottagarsidan och rapporterar:

\begin{console}
==> can_controller_tb (skipped: can_controller.vhd has no receive path yet;
    set CI_BUILD_ALL=1 to run it anyway)
\end{console}

Den kontrollen är en heuristik snarare än ett verkligt beroende, och den har ett felläge värt att
namnge: om din \chapref{c17:ch} är korrekt men stavar de uttrycken annorlunda än kapitlet får du det
här hoppmeddelandet när testbänken borde ha körts. Ingenting är trasigt; kör
\code{CI_BUILD_ALL=1 make build-project} för att tvinga fram varje testbänk oavsett. Kontrollen kan
inte fallera åt andra hållet, eftersom markören inte kan dyka upp innan en mottagväg gör det.
\end{aside}

% ------------------------------------------------------------------------------------------
\section{Att städa upp}
\label{app:sim:clean}

GHDL:s analyssteg skapar \file{work-obj93.cf} och liknande filer i den katalog du körde dem från.
Byggskriptet håller sig i stället med en egen \file{work/}-katalog, så de filerna skräpar inte ner
källträdet. Kör du de tre stegen för hand i \file{controller/} eller \file{bridge/}, städa med:

\begin{shell}
make clean
\end{shell}

som tar bort de genererade \file{work/}-katalogerna, eventuella kvarglömda \file{work-obj*.cf}, och
de vågforms- och exekverbara filer en handkörning lämnar efter sig.
